\chapter{Temporary (Transvenous) Cardiac Pacing}

\section*{Overview}
Temporary transvenous cardiac pacing is the placement of an electrode catheter into the right ventricle (or atrium) through a central vein to provide electrical pacing of the heart. It is used when the native conduction system fails or heart rate is dangerously slow.

\section*{Alternative Names}
Transvenous temporary pacemaker, emergency pacing, ventricular pacing wire placement

\section*{Principle}
An electrode catheter positioned against the endocardium delivers a small electrical impulse that depolarizes the myocardium, producing a paced rhythm. The pacemaker fires at a programmed rate when the intrinsic heart rate falls below a set threshold, maintaining cardiac output and preventing asystole.

\section*{History}
Temporary transvenous pacing was developed in the 1950s, shortly after the first external pacemakers, and became a standard tool for managing bradyarrhythmias in acute care. Alternative methods such as transcutaneous pacing are often used as a bridge to transvenous pacing.

\section*{Why It Is Done}
Ordered for symptomatic bradycardia, complete heart block with hemodynamic instability, and bradyarrhythmias following myocardial infarction or cardiac surgery. It is also used to overdrive pace certain tachyarrhythmias and to provide back-up pacing during high-risk procedures or temporary interruption of a permanent pacemaker.

\section*{Is This a Primary or Secondary Test or Procedure}
A primary therapeutic procedure for acute, potentially reversible bradyarrhythmias until a permanent pacemaker can be placed.

\section*{How It Is Performed}
Access is obtained through a central vein, most often the internal jugular, subclavian, or femoral, and a balloon-tipped pacing catheter is advanced under fluoroscopic or electrocardiographic guidance into the right ventricle. Pacing is initiated at a set rate and output, and capture is confirmed by a widened QRS followed by a T wave on the monitor. The catheter is secured and connected to an external pulse generator.

\section*{Preparation}
Informed consent when feasible, continuous electrocardiographic and blood pressure monitoring, and sterile technique for central venous access. Transcutaneous pacing pads may be placed as a bridge while the transvenous catheter is being inserted.

\section*{Contraindications and Precautions}
Contraindications include profound hypothermia with cold myocardial capture failure and severe hyperkalemia or drug toxicity where pacing may be ineffective; a reversible cause should be treated concurrently. Precautions include avoiding the subclavian approach in patients with coagulopathy or during thrombolysis, and confirming capture and sensing thresholds regularly.

\section*{Risks}
Arrhythmias, ventricular perforation with tamponade, infection, pneumothorax or bleeding from central venous access, and failure to capture or sense leading to inadequate rhythm support.

\section*{Aftercare and Recovery}
The patient is monitored in a critical care unit with continuous electrocardiography, and capture, sensing, and thresholds are checked daily. The temporary pacemaker is removed once the underlying rhythm stabilizes or a permanent pacemaker is implanted.

\section*{Outcome and Follow-up}
A paced rhythm with capture (pacing spike followed by a wide QRS and mechanical pulse) and appropriate sensing indicate correct function. Loss of capture, undersensing, or a low output threshold requiring escalation suggests catheter dislodgement, threshold rise, or generator failure.

\section*{Expected Results}
Stable pacing at the programmed rate with consistent capture, appropriate sensing of intrinsic beats, and hemodynamic stability.

\section*{Follow-up}
Continuous electrocardiographic monitoring, daily pacing threshold checks, and chest radiography to confirm catheter position. The patient is referred for permanent pacemaker implantation if the bradyarrhythmia persists.

\section*{When to Seek Medical Care}
Follow the procedure team's specific instructions. Seek urgent medical assessment for severe or worsening pain, uncontrolled bleeding, fever or signs of infection, trouble breathing, chest pain, fainting, new weakness or confusion, inability to urinate, or any rapidly worsening symptom. The appropriate warning signs depend on the procedure and the patient's medical condition.

\section*{References}
\begin{enumerate}
  \item MedlinePlus. Pacemaker placement. U.S. National Library of Medicine; 2025.
  \item Current Medical Diagnosis and Treatment. McGraw-Hill; 2025.
\end{enumerate}

\clearpage
